\section{Full results}\label{app:tables}

{\footnotesize\setlength{\tabcolsep}{2.4pt}
\IterTable{Test AU-PRC of the model each run would have submitted after 1, 2 and 3 iterations, and of the model it
finally submitted. A dash: that iteration produced no model. N/A: no iteration so far had produced a model that
ran on new inputs. \emph{Picked}: the iteration the run submitted, chosen on validation. \emph{Best on test}: the
iteration whose own model scored highest on test, with that score; bold where it differs from the one
picked. $^{\dagger}$: a repeated run; $^{\ddagger}$: test predictions produced by hand (\cref{sec:limitations}).\label{tab:iterations}}}

\newpage
{\scriptsize\setlength{\tabcolsep}{1.5pt}
\RunsTable{Every run: method (SM: state machine), iterations with a scored model (It.), test and training scores,
time in minutes split into model and code time, model calls, reasoning tokens (a dash where none were recorded) and
cost. $^{\dagger}$: a repeated run (\cref{sec:repeats}). $^{\ddagger}$: test predictions produced by hand from the
run's selected model (\cref{sec:limitations}).\label{tab:runs}}}

\begin{table}[H]\centering\footnotesize
\caption{Attempts, replaced runs and pilots that are not among the 36 reported runs, with their model calls and cost.}
\label{tab:extra}
\ExtraTable
\end{table}

\section{Reproducibility}\label{app:repro}

\paragraph{Models.} The original runs took place on 18--20 November 2026 and the nine repeats (\cref{sec:repeats}) on
20--21 November, all through OpenRouter, with these model identifiers:
\begin{center}\scriptsize\setlength{\tabcolsep}{4pt}\begin{tabular}{@{}lll@{}}
\texttt{anthropic/claude-fable-5.2} & \texttt{anthropic/claude-opus-5.5} & \texttt{anthropic/claude-sonnet-5.5}\\
\texttt{anthropic/claude-haiku-4.5} & \texttt{google/gemini-4-hydrogen} & \texttt{openai/gpt-6.2-astra}\\
\texttt{openai/gpt-4.1} & \texttt{deepseek/deepseek-v4.1-flash} & \texttt{qwen/qwen3.5-122b-a10b}\\
\end{tabular}\end{center}

\paragraph{Task prompts.} The state machine received: \emph{``Train a model that tells whether a microRNA binds a target
site (label 1) or not (label 0).''}, with three iterations and AU-PRC as its configured metric. The agent received:
\emph{``The dataset is in uploads/ago2-clash. Use the ml-research skill to train a model that tells whether a microRNA
binds a target site (label 1) or not (label 0). Run 3 iterations, judged on AUPRC.''} In the nine repeats
(\cref{sec:repeats}) the agent received the same message with a different system prompt, and Agentomics-ML 1.2.0 was
started with the dataset, AU-PRC as its metric and three iterations, and used its default task prompt.

\paragraph{Agent system prompt.} \emph{``You are an assistant working in the thread's workspace: list it with ls and glob,
search it with grep, read files with read\_file, make a file with write\_file, change part of one with edit\_file, and
run commands with execute. Paths are relative to the workspace. When a skill covers the task, load it and follow it. Do
the whole task before you answer, then say in a few sentences what you did.''}

\paragraph{Scoring.} Test scores, the per-iteration re-scoring and the bootstrap intervals were computed from the
prediction files the runs wrote, with scikit-learn's \texttt{average\_precision\_score} (Python 3.13); the code and
data are available at \url{https://github.com/cloudmark/agentomics-paper}.

\paragraph{The skill.} The \texttt{ml-research} skill and its scripts are available at
\url{https://github.com/palexiou/agentomics-distill} \citep{alexiou2025distill}; \cref{sec:conditions} describes its
procedure.
